\section{Empirical application}

\subsection{Data description}

In our analysis, we used second and third waves (NLFS 2008/09 and NLFS 2017/18) of \textit{Nepal Labor Force Survey} administrated by the National Statistics Office (NSO)(formerly known as Central Bureau of Statistics(CBS)) in collaboration with International Labor Organization (ILO) Nepal. Using a multi-stage clustering sampling approach, these survey collects rich information about individual labor pertaining to demographic profiles, labor market indicators, employment and unemployment and migration. NLFS 2008/09 (referred to as NLFS II hereafter) included interviews with 16,000 households, while the NLFS 2017/18 (referred to as NLFS III) expanded to 18,000 households. We restrict our analysis to individuals of 15-65 years who are in full time wage employment defined as individuals who work for 40 hours or more during the week. We did not include agricultural job sector in the sample, in accordance with This resulted in working-age samples of 45,222 individuals in the second round and 47,899 individuals in the third. From this, we classified employed samples comprising 6718 individuals (35\% formal and 65\% informal employment) and 6827 individuals (23\% formal and 77\% informal employment) in the NLFS II and III respectively. 

We use the measure of job-based concept of informal employment within the working-age samples. For this, we consider the non-agricultural sector . We identify employees in the formal employment as those who work in government service, financial public corporation, non-financial public corporation, governmental and non-governmental organisations and private registered financial and non-financial companies, have a permanent employment and their employer pays social security contribution or they receive the benefits of paid leave or get compensation for unused leave. Based on the same, we identify informal employment by exclusion, where we consider those not falling into these category as informal employment. 

The dependent variable in our study is the logarithm of the hourly income. The survey collects information on the usual hours work per week but not the number of weeks worked during a month. We compute hourly income by dividing monthly income by total imputed hourly work per month. Therefore, we compute the monthly hours of work by multiplying the usual hours of work per week to four. The monthly income is net income paid to people. Thus it includes transfers such as pension, social insurance deductions, and taxes as well as extra income like bonuses, overtime. premiums etc.

We investigate the determinants of wage in our study through the lens of examining human capital variables, personal and household, job and employer characteristics. We focus on experience, experience squared, and education as human capital variables, as extensive literature highlights their significant association with wage levels (Becker, 1975; Borjas, 2005; Rosen, 1972; Mincer, 1974), including the number of years with an employer (Burdett and Cole, 2010). Personal and household characteristics such as gender, caste group, marital status, and the number of children are included, given their proven impact on wages (Bardasi and Taylor, 2008; Killewald and Gough, 2013; Becker, 1975). Job characteristics encompass occupation class, vocational training, migration for work, overtime, and total household chores, all identified as influential in previous research (Becker, 1975; Borjas, 2005). Employer characteristics are captured through the job sector by economic sectors, reflecting their significant role in wage determination (Becker, 1975; Borjas, 2005).


\begin{figure}[tb] 
	\centering
	\includegraphics[width = \textwidth]{class_5_classification}
	\caption{Distribution of workers}
	\label{fig:class5Classify}
\end{figure}

Figure \ref{fig:class5Classify} illustrates the percentage distribution of employees in the formal and informal employment based on the NSCO classification. In our study, we form occupation groups for the heterogeneity based on the same, where clerical sales, agriculture trade and plant operator are recoded as clerical occupation and as it is for the elementary and managers. We justify these groups by differences in the nature of tasks, required skills levels and ssociated wages. Elementary occupation involve routine, manual tasks with lower returns to education, clerical involves moderate level of trainings and education and managers involving high-level decision-making. Figure \ref{fig:yrsschooling} shows the distribution of years of schooling for the occupation groups as mentioned above, 
 
 
 
 
 \begin{table}[htbp]
	\centering
	\footnotesize % Increase the font size
	\caption{Group characteristics by year and employment type}
	\label{tab: descriptiveoverall}
	\renewcommand{\arraystretch}{0.8} % Increase the row height
	\resizebox{\linewidth}{!}{
		\begin{tabular}{lcccc}  % Changed [lcccc] to {lcccc}
			\toprule
			\multicolumn{1}{c}{ } & \multicolumn{2}{c}{2008} & \multicolumn{2}{c}{2018} \\
			\cmidrule(l{3pt}r{3pt}){2-3} \cmidrule(l{3pt}r{3pt}){4-5}
			\text{Variables} & Informal  & Formal& Informal & Formal\\ 
			\midrule
			\text{Household size} & 5  & 4 & 5 & 4 \\
			\text{Age (in years)} & 29 & 36 & 32 & 36 \\
			\text{Experience (in years)} & 15 & 17 & 18 & 17\\
			\text{Education (in \%)} &  &  &  & \\
			\hspace{1em}\text{Illiterate} & 23 & 3.7 & 19 & 3.0\\
			\hspace{1em}\text{Below primary} & 18 & 8.0 & 18 & 4.7\\
			\hspace{1em}\text{Primary} & 27 & 14 & 29 & 13\\
			\hspace{1em}\text{Tenth grade} & 17 & 25 & 16 & 21\\
			\hspace{1em}\text{Secondary} & 7.7 & 21 & 11 & 26\\
			\hspace{1em}\text{Bachelor} & 5.2 & 18 & 5.8 & 20\\
			\hspace{1em}\text{Masters and above} & 2 & 10 & 1.9 & 12\\
			\text{Urban (in \%)} & 71 & 75 & 65 & 74\\
			\text{Married (in \%)} & 72 & 85 & 78 & 84\\
			\text{Female} & 23 & 23 & 20 & 28\\  
			\text{Hourly wage (in Rupees)} & 19  & 43  & 75  & 125 \\
			\text{No. of children} & 1  & 1  & 1  & 1 \\
			\text{Migrated for job (in \%)} & 29 & 39 & 12 & 22\\
			\text{Daily hours spent on household chores} & 0.43  & 0.57  & 0.30  & 0.56 \\
			\text{Job sector (in \%)} &  &  &  & \\
			\hspace{1em}\text{Mining\_utility} & 2.6 & 3.0 & 2.8 & 3.4\\
			\hspace{1em}\text{Construction} & 22 & 1.6 & 43 & 2.3\\
			\hspace{1em}\text{Manufacturing} & 28 & 11 & 15 & 8.1\\
			\hspace{1em}\text{Market\_services} & 26 & 20 & 21 & 24\\
			\hspace{1em}\text{Non\_Market\_services} & 13 & 62 & 16 & 60\\
			\hspace{1em}\text{Arts\_entertain} & 8.1 & 2.7 & 2.2 & 1.7\\
			\text{Occupation (in \%)} &  &  &  & \\
			\hspace{1em}\text{Elementary\_occupations} & 24 & 11 & 32 & 8.5\\
			\hspace{1em}\text{Plant\_operator} & 8.8 & 4.6 & 7.7 & 7.1\\
			\hspace{1em}\text{Agri\_trade} & 38 & 6.6 & 33 & 2.8\\
			\hspace{1em}\text{Clerical\_sales} & 15 & 26 & 12 & 24\\
			\hspace{1em}\text{Managers} & 14 & 52 & 14 & 57\\
			\bottomrule
			\text{} & & & & \\
			\text{No. of observation} & 4339& 2379 & 5272 & 1555\\
			\bottomrule
		\end{tabular}
	}
	\footnotesize{
		\raggedleft
		\textsuperscript {a} Notes: Author's calculations based on microdata from the 2008 and 2018 NLFS survey. 
	}
\end{table}

 
 
 
 
 
 \begin{figure}[tb] 
 	\centering
 	\includegraphics[width=\textwidth]{log_hourly_density_wages}
 	\caption{Log hourly density wages}
 	\label{fig:loghourlydensitywages}
 \end{figure}
 
            
   
Figure \ref{fig:loghourlydensitywages} illustrates the kernel density estimates of hourly formal and informal earnings for the elementary, clerical and managers occupation. The respective colored dashed line shows the mean wage of each of the occupations, along with the kernel density. Wage disparities within occupation group and within sectors are clealy visible, as mean wage and wage distribution is shifted to the right. The distribution of formal and informal employment log wages and their normalized earning gap at each quantile of their respective distributions is plotted in figure \ref{fig:rawwage} and figure \ref{fig:wagegap}. These figures were generated by calculating the weighted mean log wages for formal and informal employment across different quantiles for the years 2008 and 2018. 

 The figure \ref{fig:wagegap} illustrates the difference in log wages between formal and informal employment, highlighting the disparity in earnings across the wage distribution. The dashed line in the figure \ref{fig:wagegap} represents the average wage gap, providing a benchmark for the overall disparity between the two employment types. In the figure \ref{fig:rawwage} and \ref{fig:wagegap}, the wage differentials between the sectors is positive throughout the entire wage distribution for both year, with a large gap at the 25th to 60th quantile for 2008, and 35th to 70th quantile for the year 2018. 

 Table \ref{tab:overall_descriptive} displays the demographic, household, educational and job characteristics of formal and informal employment for the year 2008 and 2018 used in the earnings and selection equations for the whole, whereas table  displays these characteristics seperately for elementary, clerical and manager occupations respectively.  In terms of education, we can see that informal employment has a higher percentages of individuals with less than tenth grade education (85 \%) in 2008 and (82\%) in 2018, whereas the formal employment has 49\% in 2008 and 58\%  in 2018 for people with secondary education and above. The education levels across different occupation groups show distinct differences over the years and between employment types. In 2008, the Clerical group had 91\% of employees below the 10th grade in informal and 64.3\% in formal employment, slightly improving to 88\% and 66\% by 2018. The Elementary group consistently had the highest percentages, with 99.7\% in informal and 98\% in formal employment in 2008, decreasing to 98\% and 94.7\% in 2018. Managers had the lowest percentages, with 37.2\% below the 10th grade in informal and 31.2\% in formal employment in 2008, which significantly improved to 19.4\% and 19\% in 2018. 

Informal employees are on average younger than formal employees across the occupation class, with manager and elementary groups showing dictinct age differences. Managers in Informal employment have a distinct number of female employees with 40\% in 2008 and 49\% in 2018, which is higher than average in comparison to elementary and clerial workers. Since, experience is derived from age, they both follow the similar distribution across all the groups. Elementary employees are less likely to be in urban areas in compared to rest of the groups, whereas, informal employees in 2018 are less likely to be in urban areas compared to rest of the groups. Employees who migrated for job are consistently higher in the year 2008 across all the groups and the higher percentage of employees who have migrated have been involved in formal employment. Interestingly, managers have higher total chores hour followed by elementary workers and clerical workers. Managers are also the only group to  have the household chores hour increased from year 2008 to 2018. 



%The 17th International Conference of Labour Statisticians (ICLS) held by the International Labour Organization (ILO) in 2003 acknowledged the existence of informal employment within both formal and informal enterprises, noting that the traditional enterprise-based definition overlooked significant aspects of informality. Consequently, the 17th ICLS introduced a conceptual framework that aligned the enterprise-based concept of employment in the informal sector with a job-based concept of informal employment, thereby broadening the scope. Informal employment was defined as jobs not subject to standard labour legislation, taxation, social or legal protections to certain employment benefits regardless of whether employed by formal or informal sector enterprises.% 

\subsection{Labor market characteristics}

Since the 1990s, Nepal's labor market has transitioned from a predominantly agrarian economy to manufacturing and services based occupying both formal and informal nature of employment despite the country's contemporary socio-economic and political landscape. Economic liberalization in the 1990s, supported by IMF and World Bank reforms, diversified the economy and spurred growth in manufacturing and services, though small-scale industries showed more resilience than large-scale ones (Islam, 2014). The period of 1996 - 2006 saw the onset of a decade-long civil conflict against the backdrop of significant rural-urban income disparities, unequal and insufficient social infrastructure investment, driving widespread internal and external migration (Deraniyagala, 2006; ). Due to inadequate social safety nets, ongoing conflict, and the underdevelopment of formal institutions, the labor market faced an oversupply of workers and a shortage of formal employment opportunities. This imbalance pushed workers into the informal employment where low productivity and substandard working conditions persists.

Entering the 2000s, Nepal's large-scale manufacturing industries managed to sustain positive growth, despite experiencing fluctuations. Conversely, small-scale industries, initially resilient, began to decline significantly, facing more severe economic challenges. As a result, manufacturing  sector have emerged, where large-scale manufacturers subject to regulations and financial backings were operating formally whereas small-scale manufacturers, despite higher share of employees were operating informally. During this decade, Nepal saw infrastructural investments, particularly in strategic and rural roads, which significantly contributed to the growth of the construction sector (World Bank, 2013). The growth of urbanization and infrastructure development after escaping a decade long civil war spurred employment opportunities within construction sector, the nature of which is 

Over the years, formal employment trends have varied significantly across different sectors. The education sector, which has long been dominated by formal employment, continued to grow, reflecting the government's commitment to enhancing educational infrastructure and quality. Similarly, the finance and insurance sectors, traditionally less formalized, saw considerable increases in formal employment due to regulatory reforms and improved financial accessibility. In contrast, sectors such as manufacturing and construction faced economic hurdles that led to either stagnant or declining formal employment levels. These trends highlight the difficulties faced by these industries in maintaining formal job creation amidst global competition and fluctuating investment levels.

The informal sector remains a crucial component of Nepal’s labor market, showcasing both resilience and vulnerability. Industries such as arts, entertainment, and food \& accommodation inherently involve informal employment due to the nature of the work. These sectors saw increases in informal employment, reflecting their continued operation outside formal structures despite broader economic changes. However, the manufacturing sector's decline in both formal and informal employment suggests a contraction influenced by economic pressures and industrial shifts. Additionally, the construction sector experienced a reduction in informal employment, possibly indicating a trend towards gradual formalization or a slowdown in unregulated, small-scale activities.
